\chapter{Literature Review}
\label{ch:literature}

\section{Overview of Sign Language Recognition}
\label{sec:lit_slr_overview}
Sign Language Recognition (SLR) research has evolved from intrusive sensor-based systems (data gloves and magnetic trackers) toward vision-based recognition using standard monocular 2D RGB video cameras \cite{starner1998real, rastgoo2021sign, adaloglou2022comprehensive}. Vision-based SLR is broadly divided into Continuous Sign Language Recognition (CSLR), which transcribes unsegmented sentences \cite{koller2019deep, camgoz2018neural}, and Isolated Sign Language Recognition (ISLR), which classifies pre-segmented video clips into discrete lexical vocabulary classes. SignBridge focuses on ISLR across 76 dynamic ISL words.

\section{Convolutional Neural Networks and Transfer Learning}
\label{sec:lit_cnn_transfer}
Deep Convolutional Neural Networks (CNNs) have replaced handcrafted visual descriptors with learned spatial feature hierarchies \cite{simonyan2014very, he2016deep}. For dynamic video, Karpathy et al.~\cite{karpathy2014large} demonstrated that temporal probability averaging over individual 2D CNN frame representations provides competitive accuracy while maintaining low computational overhead. 

However, training deep CNNs from scratch requires extensive labeled data; on small gesture datasets, networks suffer severe over-parameterization and overfitting \cite{yosinski2014transferable}. Transfer learning resolves this by utilizing models pre-trained on large-scale benchmarks like ImageNet (1.2M images) \cite{deng2009imagenet}. Sandler et al.~\cite{sandler2018mobilenetv2} developed \textbf{MobileNetV2}, which leverages inverted residual blocks and linear bottlenecks to enable efficient, low-latency visual inference on resource-constrained platforms. Freezing the pre-trained ImageNet backbone preserves generalized visual primitives while adapting the dense classification head to target sign classes.

\section{Spatio-Temporal Sequence Modelling}
\label{sec:lit_temporal}
To model dynamic motion trajectories, two dominant deep learning paradigms exist: 3D CNNs (e.g., C3D \cite{tran2015learning}, I3D \cite{carreira2017quo}), which convolve across space and time simultaneously but require massive video pre-training, and recurrent hybrid architectures (CNN-LSTM) \cite{donahue2015long}. In a CNN-LSTM, 2D convolutional layers extract per-frame visual feature vectors, which an LSTM recurrent network \cite{hochreiter1997long} integrates chronologically. While theoretically appealing, training CNN-LSTMs from scratch on limited video clips faces an acute dual optimization hurdle: spatial feature extractors and recurrent transition weights must be learned simultaneously over noisy gradients \cite{adaloglou2022comprehensive}.

\section{Indian Sign Language Research and Comparative Positioning}
\label{sec:lit_isl}
Prior ISL recognition literature predominantly addresses static alphabetic finger-spelling or small vocabularies ($<35$ words) under controlled backgrounds \cite{raheja2016real, kumar2018indian}. Furthermore, existing studies rarely evaluate multiple deep learning paradigms under identical experimental splits and lack full-stack web serving systems. Table \ref{tab:lit_comparison} contrasts published paradigms with the SignBridge implementation.

\begin{table}[htbp]
\centering
\caption{Comparison of published SLR paradigms versus SignBridge.}
\label{tab:lit_comparison}
\scriptsize
\begin{tabularx}{\textwidth}{l p{3.0cm} p{2.4cm} p{3.2cm} X}
\toprule
\textbf{Work} & \textbf{Modality} & \textbf{Vocabulary} & \textbf{Architecture} & \textbf{Primary Limitation / Difference} \\
\midrule
Starner et al.~\cite{starner1998real} & Color camera & 40 ASL signs & HMM & Required colored gloves; manual features. \\
Karpathy et al.~\cite{karpathy2014large} & RGB Video & Action classes & 2D CNN + Late Fusion & General actions; not fine-grained sign gestures. \\
Donahue et al.~\cite{donahue2015long} & Video sequences & Activity videos & LRCN (CNN-LSTM) & Required large-scale video pre-training. \\
Kumar et al.~\cite{kumar2018indian} & RGB Video & ISL ($<35$ classes) & Standard 2D CNN & Limited vocabulary; no multi-model comparison. \\
Koller et al.~\cite{koller2019deep} & RGB Video & Continuous sign & Hybrid CNN-HMM & Continuous sentences; high computational cost. \\
\textbf{SignBridge} & \textbf{Standard RGB Video} & \textbf{76 ISL words} & \textbf{MobileNetV2, CNN, LSTM, Ensemble} & \textbf{76-class dynamic ISL; transfer learning vs. scratch benchmark; full FastAPI + React 19 app.} \\
\bottomrule
\end{tabularx}
\end{table}
